%%
%% Automatically generated file from DocOnce source
%% (https://github.com/hplgit/doconce/)
%%
%%


%-------------------- begin preamble ----------------------

\documentclass[%
oneside,                 % oneside: electronic viewing, twoside: printing
final,                   % draft: marks overfull hboxes, figures with paths
10pt]{article}

\listfiles               %  print all files needed to compile this document

\usepackage{relsize,makeidx,color,setspace,amsmath,amsfonts,amssymb}
\usepackage[table]{xcolor}
\usepackage{bm,ltablex,microtype}

\usepackage[pdftex]{graphicx}

\usepackage[T1]{fontenc}
%\usepackage[latin1]{inputenc}
\usepackage{ucs}
\usepackage[utf8x]{inputenc}

\usepackage{lmodern}         % Latin Modern fonts derived from Computer Modern

% Hyperlinks in PDF:
\definecolor{linkcolor}{rgb}{0,0,0.4}
\usepackage{hyperref}
\hypersetup{
    breaklinks=true,
    colorlinks=true,
    linkcolor=linkcolor,
    urlcolor=linkcolor,
    citecolor=black,
    filecolor=black,
    %filecolor=blue,
    pdfmenubar=true,
    pdftoolbar=true,
    bookmarksdepth=3   % Uncomment (and tweak) for PDF bookmarks with more levels than the TOC
    }
%\hyperbaseurl{}   % hyperlinks are relative to this root

\setcounter{tocdepth}{2}  % levels in table of contents

% Tricks for having figures close to where they are defined:
% 1. define less restrictive rules for where to put figures
\setcounter{topnumber}{2}
\setcounter{bottomnumber}{2}
\setcounter{totalnumber}{4}
\renewcommand{\topfraction}{0.95}
\renewcommand{\bottomfraction}{0.95}
\renewcommand{\textfraction}{0}
\renewcommand{\floatpagefraction}{0.75}
% floatpagefraction must always be less than topfraction!
% 2. ensure all figures are flushed before next section
\usepackage[section]{placeins}
% 3. enable begin{figure}[H] (often leads to ugly pagebreaks)
%\usepackage{float}\restylefloat{figure}

% --- fancyhdr package for fancy headers ---
\usepackage{fancyhdr}
\fancyhf{} % sets both header and footer to nothing
\renewcommand{\headrulewidth}{0pt}
\fancyfoot[LE,RO]{\thepage}
% Ensure copyright on titlepage (article style) and chapter pages (book style)
\fancypagestyle{plain}{
  \fancyhf{}
  \fancyfoot[C]{{\footnotesize \copyright\ 2013-2016, Morten Hjorth-Jensen. Released under CC Attribution-NonCommercial 4.0 license}}
%  \renewcommand{\footrulewidth}{0mm}
  \renewcommand{\headrulewidth}{0mm}
}
% Ensure copyright on titlepages with \thispagestyle{empty}
\fancypagestyle{empty}{
  \fancyhf{}
  \fancyfoot[C]{{\footnotesize \copyright\ 2013-2016, Morten Hjorth-Jensen. Released under CC Attribution-NonCommercial 4.0 license}}
  \renewcommand{\footrulewidth}{0mm}
  \renewcommand{\headrulewidth}{0mm}
}

\pagestyle{fancy}


\usepackage[framemethod=TikZ]{mdframed}

% --- begin definitions of admonition environments ---

% --- end of definitions of admonition environments ---

% prevent orhpans and widows
\clubpenalty = 10000
\widowpenalty = 10000

% --- end of standard preamble for documents ---


% insert custom LaTeX commands...

\raggedbottom
\makeindex
\usepackage[totoc]{idxlayout}   % for index in the toc
\usepackage[nottoc]{tocbibind}  % for references/bibliography in the toc

%-------------------- end preamble ----------------------

\begin{document}

% matching end for #ifdef PREAMBLE

\newcommand{\exercisesection}[1]{\subsection*{#1}}


% ------------------- main content ----------------------

% Slides for PHY981


% ----------------- title -------------------------

\thispagestyle{empty}

\begin{center}
{\LARGE\bf
\begin{spacing}{1.25}
Electromagnetic transitions and $\beta$-decay
\end{spacing}
}
\end{center}

% ----------------- author(s) -------------------------

\begin{center}
{\bf Morten Hjorth-Jensen}
\end{center}

    \begin{center}
% List of all institutions:
\centerline{{\small \href{{http://www.nscl.msu.edu/}}{National Superconducting Cyclotron Laboratory} and \href{{https://www.pa.msu.edu/}}{Department of Physics and Astronomy}, \href{{http://www.msu.edu/}}{Michigan State University}, East Lansing, MI 48824, USA}}
\end{center}
    
% ----------------- end author(s) -------------------------

% --- begin date ---
\begin{center}
Spring 2016
\end{center}
% --- end date ---

\vspace{1cm}


% !split
\subsection*{Reminder on the Wigner-Eckart theorem}

% --- begin paragraph admon ---
\paragraph{}
In our angular momentum lectures on the Wigner-Eckart theorem we developed two equations. 
One for the general expectation value that depends also on the magnetic quantum numbers  
\[
\langle \Phi^J_M|T^{\lambda}_{\mu}|\Phi^{J'}_{M'}\rangle\equiv(-1)^{J-M}\left(\begin{array}{ccc}  J & \lambda & J' \\ -M & \mu & M'\end{array}\right)\langle \Phi^J||T^{\lambda}||\Phi^{J'}\rangle,
\]
and one for the reduced matrix elements in terms of 
\[
\langle \Phi^J||T^{\lambda}||\Phi^{J'}\rangle\equiv\sum_{M,\mu,M'}(-1)^{J-M}\left(\begin{array}{ccc}  J & \lambda & J' \\ -M & \mu & M'\end{array}\right)\langle \Phi^J_M|T^{\lambda}_{\mu}|\Phi^{J'}_{M'}\rangle.
\]
% --- end paragraph admon ---





% !split
\subsection*{Why the Wigner-Eckart theorem?}

% --- begin paragraph admon ---
\paragraph{}
Unless we have observables which depend on the magnetic quantum numbers, the degeneracy 
given by these quantum numbers is not seen experimentally. 
The typical situation when we perform shell-model calculations is that the results depend on the magnetic quantum numbers. The reason for this is that it is easy to implement the Pauli principle for many particles when we work in what we dubbed for $m$-scheme.

A resulting state in a shell-model calculations will thus depend on the total value of $M$ defined as 
\[
M=\sum_{i=1}^{A}m_{j_i}.
\]
A shell model many-body state is given by a linear combination of Slater determinants $\vert \Phi_i\rangle$. That is, for some conserved quantum numbers $\lambda$  we have
\[
\vert Psi_{\lambda}\rangle = \sum_{i} C_i \vert \Phi_i\rangle,
\]
where the coefficients are the overlaps between the many-body basis sets $\Psi$ and $\Phi$ and are the resulting eigenvectors from the shell-model calculations.
% --- end paragraph admon ---




% !split
\subsection*{Why the Wigner-Eckart theorem, representing a shell-model state}

% --- begin paragraph admon ---
\paragraph{}
In second quantization, our ansatz for a state like the ground state is
\[
\vert \Phi_0\rangle = \left(\prod_{i\le F}\hat{a}_{i}^{\dagger}\right)|0\rangle,
\]
where the index $i$ defines different single-particle states up to the Fermi level. We have assumed that we have $N$ fermions.
A given one-particle-one-hole ($1p1h$) state can be written as
\[
|\Phi_i^a\rangle = \hat{a}_{a}^{\dagger}\hat{a}_i|\Phi_0\rangle,
\]
while a $2p2h$ state can be written as
\[
|\Phi_{ij}^{ab}\rangle = \hat{a}_{a}^{\dagger}\hat{a}_{b}^{\dagger}\hat{a}_j\hat{a}_i|\Phi_0\rangle,
\]
and a general $NpNh$ state as
\[
|\Phi_{ijk\dots}^{abc\dots}\rangle = \hat{a}_{a}^{\dagger}\hat{a}_{b}^{\dagger}\hat{a}_{c}^{\dagger}\dots\hat{a}_k\hat{a}_j\hat{a}_i|\Phi_0\rangle.
\]
% --- end paragraph admon ---



% !split
\subsection*{Representing a shell-model state}

% --- begin paragraph admon ---
\paragraph{}
A general shell-model many-body state 
\[
\vert Psi_{\gamma}\rangle = \sum_{i} C_i \vert \Phi_i\rangle,
\]
can be expanded as 
\[
|\Psi_{\gamma}\rangle=C_0|\Phi_0\rangle+\sum_{ai}C_i^a|\Phi_i^a\rangle+\sum_{abij}C_{ij}^{ab}|\Phi_{ij}^{ab}\rangle+\dots. 
\]
% --- end paragraph admon ---


% !split
\subsection*{Representing a shell-model state and one-body operators}

% --- begin paragraph admon ---
\paragraph{}
A one-body operator represented by a spherical tensor of rank $\lambda$ is given as 
\[
O_{\mu}^{\lambda} = \sum_{pq} \langle p \vert O_{\mu}^{\lambda} \vert q \rangle a^{\dagger}_pa_q,
\]
meaning that when we compute a transition amplitude 
\[
\langle \Psi_{\delta} \vert   O_{\mu}^{\lambda}  \vert \Psi_{\gamma} \rangle = \sum_{ij} C^{*}_{\delta i} C_{\gamma j}  \langle \Phi_i \vert O_{\mu}^{\lambda} \vert \Phi_j\rangle,
\]
we need to compute 
\[
\langle \Phi_i \vert O_{\mu}^{\lambda} \vert \Phi_j\rangle.
\]
% --- end paragraph admon ---



% !split
\subsection*{Rewriting the transition amplitude}

% --- begin paragraph admon ---
\paragraph{}
We want to rewrite 
\[
\langle \Phi_i \vert O_{\mu}^{\lambda} \vert \Phi_j\rangle, 
\]
in terms of the reduced matrix element only. 
Let us introduce the relevant quantum numbers for the states $\Phi_i$ and $\Phi_j$. We include only the relevant ones. We have then in $m$-scheme
\[
\langle \Phi_i \vert O_{\mu}^{\lambda} \vert \Phi_j\rangle = \sum_{pq} \langle p \vert O_{\mu}^{\lambda} \vert q \rangle \langle \Phi_{M}^{J} \vert a^{\dagger}_pa_q \vert \Phi_{M'}^{J'}\rangle.
\]
With a shell-model $m$-scheme basis it is straightforward to compute these amplitudes. 
However, as mentioned above, if we wish to related these elements to experiment, we need to use the Wigner-Eckart theorem and express the amplitudes in terms of reduced matrix elements.
% --- end paragraph admon ---



% !split
\subsection*{Rewriting the transition amplitude, first step}

% --- begin paragraph admon ---
\paragraph{}
We can rewrite the above transition amplitude using the Wigner-Eckart theorem. Our first step
is to rewrite the one-body operator in terms of reduced matrix elements. Since the operator is a spherical tensor we need that the annihilation operator is rewritten as (where $q$ represents $j_q$, $m_q$ etc)
\[
\tilde{a}_q=\left(-1\right)^{j_q-m_q}a_{j_q,m_q}.
\]
The operator 
\[
O_{\mu}^{\lambda} = \sum_{pq} \langle p \vert O_{\mu}^{\lambda} \vert q \rangle a^{\dagger}_pa_q,
\]
is rewritten using the Wigner-Eckart theorem as 
\[
O_{\mu}^{\lambda} = \sum_{pq} \langle p \vert \vert O^{\lambda}\vert \vert q \rangle (-1)^{j_p-m_p}\left(\begin{array}{ccc}  j_p & \lambda & j_q \\ -m_p & \mu & m_q\end{array}\right)a^{\dagger}_pa_q.
\]
% --- end paragraph admon ---




% !split
\subsection*{Rewriting the transition amplitude, second step}

% --- begin paragraph admon ---
\paragraph{}
We have
\[
O_{\mu}^{\lambda} = \sum_{pq} \langle p \vert \vert O^{\lambda}\vert \vert q \rangle (-1)^{j_p-m_p}\left(\begin{array}{ccc}  j_p & \lambda & j_q \\ -m_p & \mu & m_q\end{array}\right)a^{\dagger}_pa_q.
\]
We then single out the sum over $m_p$ and $m_q$ only and define the recoupled one-body part of the operator as 
\[
\lambda^{-1}\left[ a_{j_p}^{\dagger}\tilde{a}_{j_q}\right]^{\lambda}_{\mu} = \sum_{m_p,m_q} (-1)^{j_p-m_p}\left(\begin{array}{ccc}  j_p & \lambda & j_q \\ -m_p & \mu & m_q\end{array}\right)a^{\dagger}_pa_q,
\]
with $\lambda=\sqrt{2\lambda+1}$.
This gives the following expression for the one-body operator
\[
O_{\mu}^{\lambda} = \sum_{j_pj_q} \langle p \vert \vert O^{\lambda}\vert \vert q \rangle \lambda^{-1}\left[ a_{j_p}^{\dagger}\tilde{a}_{j_q}\right]^{\lambda}_{\mu}.
\]
% --- end paragraph admon ---




% !split
\subsection*{Rewriting the transition amplitude, third step}

% --- begin paragraph admon ---
\paragraph{}
With 
\[
O_{\mu}^{\lambda} = \sum_{j_pj_q} \langle p \vert \vert O^{\lambda}\vert \vert q \rangle \lambda^{-1}\left[ a_{j_p}^{\dagger}\tilde{a}_{j_q}\right]^{\lambda}_{\mu},
\]
we can write 
\[
\langle \Phi^J_M|O^{\lambda}_{\mu}|\Phi^{J'}_{M'}\rangle = \sum_{pq} \langle p \vert O_{\mu}^{\lambda} \vert q \rangle \langle \Phi_{M}^{J} \vert a^{\dagger}_pa_q \vert \Phi_{M'}^{J'}\rangle,
\]
as
\[
\langle \Phi^J_M|O^{\lambda}_{\mu}|\Phi^{J'}_{M'}\rangle = \sum_{j_pj_q} \langle p \vert \vert O^{\lambda} \vert\vert q \rangle \langle \Phi_{M}^{J}\vert \lambda^{-1}\left[ a_{j_p}^{\dagger}\tilde{a}_{j_q}\right]^{\lambda}_{\mu}  \vert \Phi_{M'}^{J'}\rangle.
\]
We have suppressed the summation over quantum numbers like $n_p,n_q$ etc.
% --- end paragraph admon ---




% !split
\subsection*{Rewriting the transition amplitude, final step}

% --- begin paragraph admon ---
\paragraph{}
Using the Wigner-Eckart theorem
\[
\langle \Phi^J_M|O^{\lambda}_{\mu}|\Phi^{J'}_{M'}\rangle\equiv(-1)^{J-M}\left(\begin{array}{ccc}  J & \lambda & J' \\ -M & \mu & M'\end{array}\right)\langle \Phi^J||O^{\lambda}||\Phi^{J'}\rangle,
\]
we can then define 
\[
\langle \Phi^J||O^{\lambda}||\Phi^{J'}\rangle = \lambda^{-1}\sum_{j_pj_q} \langle p \vert \vert O^{\lambda} \vert\vert q \rangle \langle \Phi_{M}^{J}\vert\vert \left[ a_{j_p}^{\dagger}\tilde{a}_{j_q}\right]^{\lambda}\vert  \vert \Phi_{M'}^{J'}\rangle.
\]
The quantity to the left in the last equation is normally called the transition amplitude or in case of a decay process, simply the decay amplitude.  The quantity 
$\langle \Phi_{M}^{J}\vert \lambda^{-1}\left[ a_{j_p}^{\dagger}\tilde{a}_{j_q}\right]^{\lambda}_{\mu}  \vert \Phi_{M'}^{J'}\rangle$ is called the one-body transition density while the corresponding reduced one is simply called the reduced one-body transition density. 
The transition densities characterize the many-nucleon
properties of the initial  and final states. They do not
carry information about the transition operator beyond its one-body character. 
Finally, note that in a shell-model calculation it is actually $\langle \Phi^J_M|O^{\lambda}_{\mu}|\Phi^{J'}_{M'}\rangle$ which is calculated. The exercises here are done in order to related to experimental results which do not depend on the magnetic substates.
% --- end paragraph admon ---







% !split
\subsection*{Electromagnetic multipole moments and transitions}

% --- begin paragraph admon ---
\paragraph{}
The reduced
transition probability $B$ is defined in terms of reduced matrix
elements of a one-body operator by
\[
B(i \rightarrow f)= \frac{\langle J_{f}||{\cal O}(\lambda )||J_{i}\rangle^{2}}{(2J_{i}+1)}. 
\]
With our definition of the reduced matrix element,
\[
\langle J_{f}||{\cal O}(\lambda )||J_{i}\rangle^{2} =\langle J_{i}||{\cal O}(\lambda )||J_{f}\rangle^{2},  
\]
the transition probability $B$ depends upon the direction of the transition by the factor
of $(2J_{i}+1)$. For electromagnetic
transitions $J_{i}$ is that for the higher-energy initial state. But in
Coulomb excitation the initial state is usually
taken as the ground state, and it is normal to use the notation $B(\uparrow)$ for transitions from the ground state.
% --- end paragraph admon ---



% !split
\subsection*{Electromagnetic multipole moments and transitions}

% --- begin paragraph admon ---
\paragraph{}
The one-body operators ${\cal O}(\lambda)$ represent a sum over
the operators for the individual nucleon degrees of freedom  $i$
\[
{\cal O}(\lambda ) = \sum_{i} O(\lambda ,i). 
\]

The electric transition operator is given by
\[
O(E\lambda ) = r^{\lambda } \; Y^{\lambda }_{\mu }(\hat{r}) \; e_{q} e, 
\]
were $Y^{\lambda }_{\mu }$ are the spherical harmonics
and $q$ stands for proton $q=p$ or neutron $q=n$.
% --- end paragraph admon ---




% !split
\subsection*{Electromagnetic multipole moments and transitions}

% --- begin paragraph admon ---
\paragraph{}
Gamma transitions
with $\lambda=0$ are forbidden because the photon must carry off
at least one unit of angular momentum. The $e_{q}$
are the electric charges for the proton and neutron in units of $e$.
For the free-nucleon
charge we would take $e_{p}=1$ and $e_{n}=0$, for the
proton and neutron, respectively.
Although the bare operator acts upon the protons,
we will keep the general expression in terms of $e_{q}$ in order
to incorporate the \textbf{effective charges} for the proton and
neutron, which represent the center-of-mass corrections and the
average effects of the renormalization from wavefunction
admixtures outside the model space.
% --- end paragraph admon ---



% !split
\subsection*{Electromagnetic multipole moments and transitions}

% --- begin paragraph admon ---
\paragraph{}
The magnetic transition operator is given by:
\[
O(M\lambda)=\left[\mathbf{l}\frac{2g^{l }_{q}}{(\lambda +1)}+ \mathbf{s}g^{s}_{q}\right]\mathbf{\nabla}[r^{\lambda }Y^{\lambda }_{\mu }(\hat{r})]\mu _{N}
\]
\[
= \sqrt{\lambda (2\lambda +1)}\left[[Y^{\lambda -1}(\hat{r})\otimes \mathbf{l}\,]^{\lambda }_{\mu }\frac{2g^{l}_{q}}{(\lambda +1)}
+ [Y^{\lambda -1}(\hat{r})\otimes \mathbf{s}\,]^{\lambda }_{\mu }g^{s}_{q}\right]r^{\lambda -1}\mu _{N}, 
\]
 where $\mu_{N}$ is the nuclear magneton,
\[
\mu _{N}=\frac{e\hbar }{2m_{p}c} = 0.105 \; e {\rm fm}, 
\]
and where $m_{p}$ is the mass of the proton.
% --- end paragraph admon ---



% !split
\subsection*{Electromagnetic multipole moments and transitions}

% --- begin paragraph admon ---
\paragraph{}
The g-factors $g^{l}_{q}$ and $g^{s}_{q}$
are the orbital and spin g-factors
for the proton and neutron, respectively.
The
free-nucleon values for the g-factors are $g^{l}_{p}=1$, $g^{l}_{n}=0$,
$g^{s}_{p}=5.586$ and $g^{s}_{n}=-3.826$. We may use effective values
for these g-factors to take into account the truncation of the model
space.
% --- end paragraph admon ---



% !split
\subsection*{Electromagnetic multipole moments and transitions}

% --- begin paragraph admon ---
\paragraph{}
The most common types of transitions are $E1$, $E2$ and $M1$.
The $E1$ transition operator is given by $\lambda=1$
\[
O(E1) = rY^{(1)}_{\mu }(\hat{r})e_{q} e= \sqrt{\frac{3}{4\pi }}\mathbf{r}e_{q} e.
\]
The $E2$ transition operator with $\lambda=2$
\[
O(E2) = r^{2}Y^{(2)}_{\mu }(\hat{r})e_{q} e,  
\]
% --- end paragraph admon ---



% !split
\subsection*{Electromagnetic multipole moments and transitions}

% --- begin paragraph admon ---
\paragraph{}
The $M1$ transition operator with $\lambda=1$ and with
\[
  Y^{0}=1/\sqrt{4\pi },
\]
we have
\[
O(M1)=\sqrt{\frac{3}{4\pi }}[\mathbf{l}g^{l }_{q}+\mathbf{s} \; g^{s}_{q}]\mu _{N}. 
\]
% --- end paragraph admon ---



% !split
\subsection*{Electromagnetic multipole moments and transitions}

% --- begin paragraph admon ---
\paragraph{}
The selection rules are given by the triangle condition for the
angular momenta, $\Delta(J_{i},J_{f},\lambda)$. 

The electromagnetic
interaction conserves parity, and the elements of the
operators for $E\lambda$ and $M\lambda$
can be classified according to their transformation under parity
change
\[
\hat{P}\hat{O}\hat{P}^{-1}=\pi_{O}\hat{O}, 
\]
where we have $\pi _{O}=(-1)^{\lambda }$ for $Y^{\lambda }$,
$\pi _{O}=-1$ for the vectors
$\mathbf{r}$, $\mathbf{\nabla}$ and $\mathbf{p}$, and $\pi _{O}=+1$ for the 
pseudo vectors
$\mathbf{l}=\mathbf{r}\times\mathbf{p}$ and $\mathbf{\sigma}$. For a given matrix element we have:
\[
\langle\Psi _{f}\vert {\cal O}\vert \Psi _{i}\rangle =\langle\Psi _{f}\vert P^{-1}P{\cal O}P^{-1}P\vert \Psi _{i}\rangle=\pi _{i}\pi _{f}\pi _{O} \langle\Psi _{f}\vert {\cal O}\vert \Psi _{i}\rangle. 
\]
The matrix element will vanish unless  $\pi _{i}\pi _{f}\pi _{O}=+1$.
% --- end paragraph admon ---



% !split
\subsection*{Electromagnetic multipole moments and transitions}

% --- begin paragraph admon ---
\paragraph{}
The transitions are divided into two classes, those
which do not change parity change $\pi _{i}\pi _{f}=+1$ which go by the
operators with $\pi _{O}=+1$:
\[
\pi _{i}\pi _{f}=+1 \; \mathrm{for} \; M1, E2, M3, E4 \dots,
\]
and the ones which do change parity change $\pi _{i}\pi _{f}=-1$
which go by the operators with $\pi _{O}=-1$:
\[
\pi _{i}\pi _{f}=-1 \; \mathrm{for} \; E1, M2, E3, M4 \dots.
\]
% --- end paragraph admon ---



% !split
\subsection*{Electromagnetic multipole moments and transitions}

% --- begin paragraph admon ---
\paragraph{}
The 
electromagnetic moment operator can be expressed in terms of the
electromagnetic transition operators.
By the parity selection rule of the moments are nonzero
only for $M1$, $E2$, $M3$, $E4,\ldots$.
The most common are:
\[
\mu =\sqrt{\frac{4\pi }{3}}\langle J,M=J\vert {\cal O}(M1)\vert J,M=J\rangle= \sqrt{\frac{4\pi }{3}}\left\{\begin{array}{ccc} J & 1 &J \\  -J &  0 &  J \end{array}\right\}
\langle J\vert \vert {\cal O}(M1)\vert \vert J\rangle,
\]
and
\[
Q = \sqrt{\frac{16\pi }{5}}\langle J,M=J\vert {\cal O}(E2)\vert J,M=J\rangle= \sqrt{\frac{16\pi }{5}}\,
  \left(\begin{array}{ccc}  {J}&  {2} & {J}\\  {-J} & {0}&  {J}\end{array}\right)\langle J\vert\vert {\cal O}(E2)\vert\vert J\rangle.
\]
% --- end paragraph admon ---



% !split
\subsection*{Electromagnetic multipole moments and transitions}

% --- begin paragraph admon ---
\paragraph{}
Electromagnetic transitions and moments depend upon the reduced nuclear
matrix elements $\langle f\vert\vert {\cal O}(\lambda )\vert\vert i\rangle$. These can be expressed as a sum over one-body transition
densities (OBTD) times single-particle matrix elements
\[
\langle f\vert\vert {\cal O}(\lambda )\vert\vert i\rangle
=\sum _{k_{\alpha } k_{\beta }}\mathrm{OBTD}(f i k_{\alpha } k_{\beta } \lambda )
 \langle k_{\alpha }\vert\vert O(\lambda )\vert\vert k_{\beta }\rangle, 
\]
where the OBTD is given by
\[
\mathrm{OBTD}(f i k_{\alpha} k_{\beta}\lambda)= \frac{\langle f\vert\vert [a^{+}_{k_{\alpha }}\otimes \tilde{a}_{k_{\beta }}]^{\lambda }\vert\vert i\rangle}{\sqrt{(2\lambda +1)}}. 
\]
The labels $i$ and $f$ are a short-hand notation for the initial
and final state quantum numbers $(n \omega _{i}J_{i})$ and $(n\omega_{f}J_{f})$,
respectively. Thus the problem is divided into two parts, one
involving the nuclear structure dependent one-body transition
densities OBTD, and the other involving the reduced
single-particle matrix
elements (SPME).
% --- end paragraph admon ---





% !split
\subsection*{Electromagnetic multipole moments and transitions}

% --- begin paragraph admon ---
\paragraph{}
The SPME for the $E\lambda$ operator is given by
\[
\langle k_{a}\vert\vert O(E\lambda )\vert\vert k_{b}\rangle=(-1)^{j_{a}+1/2}\frac{[1+(-1)^{l_{a}+\lambda +l_{b}}]}{2}
\]
\[
 \times\sqrt{ {(2j_{a}+1)(2\lambda +1)(2j_{b}+1)\over4\pi }}\left(\begin{array}{ccc}  {j_{a}}&  {\lambda} &  {j_{b}}\\  {1/2} & {0}&  {-1/2}\end{array}\right)\langle k_{a}\vert r^{\lambda }\vert k_{b}\rangle e_{q} e.
\]
% --- end paragraph admon ---



% !split
\subsection*{Electromagnetic multipole moments and transitions}

% --- begin paragraph admon ---
\paragraph{}
The SPME for the spin part of the magnetic operator is
\[
\langle k_{a}\vert\vert O(M\lambda ,s)\vert\vert k_{b}\rangle =
\]
\[
=\sqrt{\lambda (2\lambda +1)}<j_{a}\vert\vert [Y^{\lambda -1}(\hat{r})\otimes\mathbf{s}\,]^{\lambda }\vert\vert j_{b}><k_{a}\vert r^{\lambda -1}\vert k_{b}>g^{s}_{q}\mu _{N},
\]
\[
= \sqrt{\lambda (2\lambda +1)}\, \sqrt{(2j_{a}+1)(2j_{b}+1)(2\lambda +1)}\left\{\begin{array}{ccc}  {l _{a}}&  {1/2} & {j_{a}}\\  {l _{b}}&  {1/2} & {j_{b}}\\  {\lambda -1} & {1} & {\lambda}\end{array}\right\}
\]
\[
\times \langle l _{a}\vert\vert Y^{\lambda -1}(\hat{r})\vert\vert l _{b}\rangle\langle\vert\vert \mathbf{s}\vert\vert s\rangle\langle k_{a}\vert r^{\lambda -1}\vert k_{b}\rangle g^{s}_{q}\mu _{N},
\]
where
\[
\langle\vert\vert \mathbf{s}\vert\vert s\rangle = \sqrt{3/2}.
\]
% --- end paragraph admon ---



% !split
\subsection*{Electromagnetic multipole moments and transitions}

% --- begin paragraph admon ---
\paragraph{}
The SPME for the orbital part of the magnetic operator is:
\[
\langle k_{a}\vert\vert O(M\lambda ,l )\vert\vert k_{b}\rangle=
\]
\[
= \frac{\sqrt{\lambda (2\lambda +1)}\, }{\lambda +1}
\langle j_{a}\vert\vert [Y^{\lambda -1}(\hat{r})\otimes\mathbf{l}\,]^{\lambda }\vert\vert j_{b}\rangle
\langle k_{a}\vert r^{\lambda -1}\vert k_{b}\rangle g^{l }_{q}\mu _{N}
\]
\[
=\frac{\sqrt{\lambda (2\lambda +1)}\, }{\lambda +1}(-1)^{l _{a}+1/2+j_{b}+\lambda } \sqrt{(2j_{a}+1)(2j_{b}+1)}
\]
\[
\times\left\{\begin{array}{ccc}  {l _{a}} &  {l _{b}} & {\lambda} \\  {j_{b}}&  {j_{a}}&  {1/2}\end{array}\right\}
\langle l _{a}\vert\vert [Y^{\lambda -1}(\hat{r})\otimes\mathbf{l}\,]^{\lambda }\vert\vert l _{b}\rangle
\langle k_{a}\vert r^{\lambda -1}\vert k_{b}\rangle g^{l }_{q}\mu _{N}, 
\]
% --- end paragraph admon ---



% !split
\subsection*{Electromagnetic multipole moments and transitions}

% --- begin paragraph admon ---
\paragraph{}
where we have defined
\[
\langle l _{a}\vert\vert [Y^{\lambda -1}(\hat{r})\otimes\mathbf{l}]^{\lambda }\vert\vert l _{b}\rangle=(-1)^{\lambda +l _{a}+l _{b}} \sqrt{(2\lambda +1)l _{b}(l _{b}+1)(2l _{b}+1)}
\]
\[
\times\left\{\begin{array}{ccc}  {\lambda -1} & {1}&  {\lambda}\\   {l _{b}}&  {l _{a}} &  {l _{b}} \end{array}\right\}
\langle l _{a}\vert\vert Y^{\lambda -1}(\hat{r})\vert\vert l _{b}\rangle,
\]
with
\[
\langle l _{a}\vert\vert Y^{\lambda -1}(\hat{r})\vert\vert l _{b}\rangle=(-1)^{l _{a}} \sqrt{\frac{(2l _{a}+1)(2l _{b}+1)(2\lambda -1)}{4\pi }}\left(\begin{array}{ccc}  {l _{a}} & {\lambda -1} & {l _{b}}\\  {0} & {0}&  {0}\end{array}\right)
. 
\]
% --- end paragraph admon ---



% !split
\subsection*{Electromagnetic multipole moments and transitions}

% --- begin paragraph admon ---
\paragraph{}
For the $M1$ operator the radial matrix element is
\[
<k_{a}\vert r^{0}\vert k_{b}>\, = \delta _{n_{a},n_{b}},
\]
and the SPME simplify to:
\[
\langle k_{a}\vert\vert O(M1,s)\vert\vert k_{b}\rangle=\sqrt{ \frac{3}{4\pi }}\langle j_{a}\vert\vert \mathbf{s}\,\vert\vert j_{b}\rangle \delta _{n_{a},n_{b}}g^{s}_{q}\mu_{N}
\]
\[
=\sqrt{ \frac{3}{4\pi }}(-1)^{l _{a}+j_{a}+3/2}
\sqrt{(2j_{a}+1)(2j_{b}+1)}\left\{\begin{array}{ccc}  {1/2}&  {1/2} & {1} \\ {j_{b}} & {j_{a}}&  {l _{a}}\end{array}\right\}
\]
\[
\times\langle s\vert\vert \mathbf{s}\vert\vert s\rangle \delta _{l _{a},l _{b}} \delta _{n_{a},n_{b}}g^{s}_{q}\mu _{N},
\]
% --- end paragraph admon ---



% !split
\subsection*{Electromagnetic multipole moments and transitions}

% --- begin paragraph admon ---
\paragraph{}
where we have
\[
<s\vert\vert \mathbf{s}\,\vert\vert s>\, = \sqrt{3/2}\, ,
\]
and
\[
<k_{a}\vert\vert O(M1,l )\vert\vert k_{b}>\, =
\sqrt{ \frac{3}{4\pi }}\,
<j_{a}\vert\vert \mathbf{l}\,\vert\vert j_{b}> \delta _{n_{a},n_{b}} \; g^{l }_{q} \; \mu _{N}
\]
\[
=\sqrt{ \frac{3}{4\pi }}\,
(-1)^{l _{a}+j_{b}+3/2} \sqrt{(2j_{a}+1)(2j_{b}+1)}\,
   \left\{\begin{array}{ccc}  {l _{a}}&  {l _{b}} & {1}\\  {j_{b}}&  {j_{a}}&  {1/2}\end{array}\right\}
\]
\[
\times\langle l _{a}\vert\vert \mathbf{l}\,\vert\vert l _{b}\rangle \delta _{n_{a},n_{b}}g^{l }_{q}\mu _{N} ,       
\]
where
\[
\langle l _{a}\vert\vert \mathbf{l}\,\vert\vert l _{b}\rangle = \delta _{l _{a},l _{b}} \sqrt{l _{a}(l_{a}+1)(2l _{a}+1)}.
\]
Thus the $M1$ operator connects only
those orbitals which have the same $n$ and $l$ values.
% --- end paragraph admon ---




% !split
\subsection*{$\beta$-decay}

% --- begin paragraph admon ---
\paragraph{}
We will now focus on  allowed $\beta$-decay.  
Suhonen's chapter 7 and Alex Brown's chapter 29 cover much of the material to be discussed on $\beta$-decay.
% --- end paragraph admon ---



% !split
\subsection*{$\beta$-decay}

% --- begin paragraph admon ---
\paragraph{}
The allowed beta decay rate $W$ between a specific set of
initial and final states is given by
\[
W_{i,f} = (f/K_{o}) \left[ g_{V}^{2} \; B_{i,f}(F_{\pm})+ g_{A}^{2}B_{i,f}(GT_{\pm})\right], 
\]
where $f$ is dimensionless three-body
phase-space factor which depends upon the
beta-decay $Q$ value,
and $K_{o}$ is a specific combination of fundamental constants
\[
  K_{o}=\frac{2\pi^{3}\hbar^{7}}{ m_{e}^{5} c^{4}}= 1.8844 \times 10^{-94}\mathrm{erg}^{2}\mathrm{cm}^{6}\mathrm{s}. 
\]
The $\pm$ signrefer to $\beta_{\pm}$ decay of nucleus
$(A_{i},Z_{i})$ into nucleus $(A_{i},Z_{i} \mp 1)  $.
The weak-interaction vector ($V$) and axial-vector ($A$) coupling
constants for the decay of neutron into a proton are denoted by $g_{V}$
and $g_{A}$, respectively.
% --- end paragraph admon ---



% !split
\subsection*{$\beta$-decay}

% --- begin paragraph admon ---
\paragraph{}
The total decay rate for a given
initial state is obtained by summing the partial rates over all
final states
\[
W = \displaystyle\sum _{f} W_{if}, 
\]
with the branching fraction to a specific final state given by
\[
b_{if} = \frac{W_{if}}{W}.  
\]
Beta decay lifetime are usually given in terms of the half-life with
a total half-life of
\[
T_{1/2} = \frac{{\rm ln}(2)}{W}.
\]
The partial half-life for a particular final state will be
denoted by $t_{1/2}$
\[
t_{1/2} = \frac{T_{1/2}}{b_{if}}. 
\]
% --- end paragraph admon ---



% !split
\subsection*{$\beta$-decay}

% --- begin paragraph admon ---
\paragraph{}
Historically
one combines the partial half-life for a particular decay
with the calculated
phase-space factor $f$ to obtain  an \textbf{ft} value given by
\[
  f t_{1/2}=\frac{C }{\left[B(F_{\pm})+(g_{A}/g_{V})^{2} B(GT_{\pm}) \right] }  
\]
where
\[
  C = \frac{{\rm ln}(2)K_{o}}{(g_{V})^{2}}. 
\]
% --- end paragraph admon ---



% !split
\subsection*{$\beta$-decay}

% --- begin paragraph admon ---
\paragraph{}
One often compiles the allowed beta decay in terms of a \textbf{logft}
which stands for log$_{10}$ of the $ft_{1/2}$ value.

     The values of the coupling constants for Fermi decay,
$g_{V}$, and Gamow-Teller decay, $g_{A}$ are obtained as follows.
For a $0^{+} \rightarrow  0^{+}$ nuclear transition $B(GT)=0$, and for a
transition between $T=1$ analogue states with $B(F)=2$ we find 
\[
       C = 2 t_{1/2} f.  
\]
The partial half-lives and $Q$ values for several $0^{+} \rightarrow  0^{+}$ analogue
transitions have been measured to an accuracy of about one part in
10000. With  phase space factors one obtains 
\[
      C = 6170(4)    
\]
This result, together with the value of $K_{o}$ can be used to obtain $g_{V}$.
% --- end paragraph admon ---



% !split
\subsection*{$\beta$-decay}

% --- begin paragraph admon ---
\paragraph{}
At the quark level $g_{V}=-g_{A}$.
But for nuclear structure we use the value obtained from the
neutron to proton beta decay
\[
     \vert g_{A}/g_{V}\vert  = 1.261(8). 
\]
% --- end paragraph admon ---



% !split
\subsection*{$\beta$-decay}

% --- begin paragraph admon ---
\paragraph{}
The operator for Fermi beta decay in terms of sums
over the nucleons is
\[
{\cal O}(F_{\pm}) =  \displaystyle\sum _{k} t_{k\pm}.
\]
The matrix element is
\[
B(F) =\, \vert \langle f\vert  T_{\pm} \vert i\rangle\vert ^{2}, 
\]
where
\[
T_{\pm} = \sum _{k} t_{\pm}  
\]
is the total isospin raising and lowering operator for total
isospin constructed out of the
basic nucleon isospin raising and lowering operators
\[
t_{-}\vert n\rangle=\vert p\rangle \hspace{1cm} t_{-}\vert p\rangle= 0,
\]
and
\[
t_{+}\vert p\rangle=\vert n\rangle, \;\; t_{+}\vert n\rangle= 0.  
\]
% --- end paragraph admon ---



% !split
\subsection*{$\beta$-decay}

% --- begin paragraph admon ---
\paragraph{}
The matrix elements obey the triangle
conditions $J_{f}=J_{i}$ ($\Delta J=0$). The Fermi operator has
$\pi _{O}=+1$, and thus the
initial and final nuclear states must have $\pi _{i}\pi _{f}=+1$ for
the matrix element to be
nonzero under the parity transform.

When isospin is conserved
the Fermi matrix element must obey the isospin triangle condition
$T_{f}=T_{i}$ $(\Delta T=0)$, and the Fermi operator can only connect
isobaric analogue states.
% --- end paragraph admon ---



% !split
\subsection*{$\beta$-decay}

% --- begin paragraph admon ---
\paragraph{}
For $\beta_{-}$ decay
\[
T_{-}\vert \omega _{i},J_{i},M_{i},T_{i},T_{zi}\rangle
\]
\[
= \sqrt{(T_{i}(T_{i}+1)-T_{zi}(T_{zi}-1)}\vert \omega _{i},J_{i},M_{i},T_{i},T_{zi}-1\rangle,      
\]
and
\[
B(F_{-}) =\, \vert \langle \omega _{f},J_{f},M_{f},T_{f},T_{zi}-1\vert T_{-}\vert \omega 
_{i},J_{i},M_{i},T_{i},T_{zi}\rangle\vert ^{2}
\]
\[
= [T_{i}(T_{i}+1)-T_{zi}(T_{zi}-1)]
\delta _{\omega _{f},\omega }\;\delta _{J_{i},J_{f}}\;\delta _{M_{i},M_{f}}\;\delta _{T_{i},T_{f}}. 
\]
% --- end paragraph admon ---



% !split
\subsection*{$\beta$-decay}

% --- begin paragraph admon ---
\paragraph{}
For $\beta_{ + }$ we have
\[
B(F_{+}) =\, \vert \langle \omega _{f},J_{f},M_{f},T_{f},T_{zi}+1\vert T_{+}\vert \omega 
_{i},J_{i},M_{i},T_{i},T_{zi}\rangle\vert ^{2}
\]
\[
= [T_{i}(T_{i}+1)-T_{zi}(T_{zi}+1)]
\delta _{\omega _{f},\omega }\;\delta _{J_{i},J_{f}}\delta _{M_{i},M_{f}}\delta _{T_{i},T_{f}}. 
\]
For neutron-rich nuclei ($N_{i}> Z_{i}$) we have $T_{i}=T_{zi}$ and thus
\[
B(F_{-})(N_{i}> Z_{i}) = 2T_{zi} = (N_{i}-Z_{i})\delta _{\omega _{f},\omega }\delta _{J_{i},J_{f}}\;\delta _{M_{i},M_{f}}\;\delta _{T_{i},T_{f}}, 
\]
and
\[
B(F_{+})(N_{i}> Z_{i}) = 0. 
\]
% --- end paragraph admon ---



% !split
\subsection*{$\beta$-decay}

% --- begin paragraph admon ---
\paragraph{}
The reduced single-particle
matrix elements are given by
\[
\langle k_{a},p\vert\vert \sigma  t_{-}\vert\vert k_{b},n\rangle=\langle k_{a},n\vert\vert \sigma  t_{+}\vert\vert k_{b},p\rangle= 2\langle k_{a}\vert\vert \mathbf{s}\vert\vert k_{b}\rangle, 
\]
where the matrix elements of $\mathbf{s}$ are given by
\[
\langle k_{a}\vert\vert \mathbf{s}\vert\vert k_{b}\rangle=\langle j_{a}\vert\vert \mathbf{s}\vert\vert j_{b}\rangle \delta_{n_{a},n_{b}}
\]
\[
=(-1)^{l_{a}+j_{a}+3/2}
\sqrt{(2j_{a}+1)(2j_{b}+1)}\left\{\begin{array}{ccc}  {1/2}&  {1/2}&  {1} \\ {j_{b}}&  {j_{a}}&  {l_{a}}\end{array}\right\}
\langle s\vert\vert \mathbf{s}\,\vert\vert s\rangle \delta _{\ell _{a},\ell _{b}} \delta_{n_{a},n_{b}}
,
\]
with
\[
\langle s\vert\vert \mathbf{s}\vert\vert s\rangle= \sqrt{3/2}.
\]
% --- end paragraph admon ---



% !split
\subsection*{$\beta$-decay}

% --- begin paragraph admon ---
\paragraph{}
The matrix elements of $\mathbf{s}$ has the selection rules $\delta_{ \ell_{a} , \ell_{b} }$
and $\delta_{n _{a} ,n _{b} }$. Thus the orbits which are connected by the GT operator
are very selective; they are those in the same major oscillator shell
with the same $\ell$ value. The matrix elements such as
$1s_{1/2}-0d_{3/2}$ which have the allowed $\Delta j$ coupling but
are zero due to the $\Delta\ell$ coupling are called $\ell$-\textbf{forbidden} matrix
elements.
% --- end paragraph admon ---



% !split
\subsection*{$\beta$-decay}

% --- begin paragraph admon ---
\paragraph{}
Sum rules for Fermi and Gamow-Teller matrix elements can be obtained easily.

The sum rule for Fermi is obtained from the sum
\[
\sum _{f} \left[ B_{fi}(F_{-}) - B_{fi}(F_{+}) \right]
=\sum _{f} \left[ \vert \langle f\vert  T_{-} \vert i\rangle\vert ^{2} -  \vert \langle f\vert  T_{+} \vert 
i\rangle\vert ^{2} \right]
\]
The final states $f$ in the $T_{-}$ matrix element go
with the $Z_{f}=Z_{i}+1$ nucleus and those in the $T_{+}$ matrix element
to with the $Z_{f}=Z_{i}-1$ nucleus. One can explicitly sum over the
final states to obtain
\[
\sum _{f} \left[ \langle i\vert  T_{+} \vert f\rangle \langle f\vert T_{-}\vert i\rangle -  \langle i\vert  T_{-} \vert f\rangle \langle f\vert T_{+}\vert i\rangle \right]
\]
\[
= \langle i\vert  T_{+} T_{-} -  T_{-} T_{+}\vert i\rangle =\langle i\vert  2T_{z}\vert i\rangle  = (N_{i}-Z_{i}).       
\]
% --- end paragraph admon ---



% !split
\subsection*{$\beta$-decay}

% --- begin paragraph admon ---
\paragraph{}
The sum rule for Gamow-Teller is obtained as follows
\[
\sum_{f,\mu}  \vert \langle f\vert  \sum_{k} \sigma_{k,\mu} t_{k-} \vert i\rangle\vert^{2}
- \sum_{f,\mu} \vert \langle f\vert  \sum_{k} \sigma_{k,\mu } t_{k+} \vert i\rangle\vert^{2}
\]
\[
= \sum_{f,\mu}\langle i\vert  \sum_{k} \sigma_{k,\mu} t_{k+} \vert f\rangle\langle f\vert  \sum_{k'} 
\sigma_{k',\mu} t_{k'-} \vert i\rangle
\]
\[
-  \sum_{f,\mu}
\langle i\vert  \sum_{k} \sigma_{k,\mu } t_{k-} \vert f\rangle\langle f\vert\sum_{k'} 
\sigma_{k',\mu } t_{k'+} \vert i\rangle
\]
\[
= \sum_{\mu} \left[\langle i\vert  \left(\sum _{k} \sigma_{k,\mu} t_{k+} \right)
    \left( \sum_{k'} \sigma_{k',\mu} t_{k'-}\right)
-   \left( \sum_{k} \sigma_{k,\mu} t_{k-} \right)
\left( \sum_{k'} \sigma_{k',\mu} t_{k'+} \right) \vert i\rangle
\right]
\]
\[
= \sum_{\mu }
\langle i\vert  \sum_{k}  \sigma ^{2}_{k,\mu } \left[ t_{k+} t_{k-} - t_{k-} t_{k+} \right] 
\vert i\rangle
= 3
\langle i\vert  \sum_{k} \left[ t_{k+} t_{k-} - t_{k-} t_{k+} \right] \vert i\rangle
\]
\[
=   3\langle i\vert  T_{+} T_{-} -  T_{-} T_{+}\vert i\rangle= 3\langle i\vert  2T_{z}\vert i\rangle  = 3(N_{i}-Z_{i}). 
\]
% --- end paragraph admon ---



% !split
\subsection*{$\beta$-decay}

% --- begin paragraph admon ---
\paragraph{}
We have used the fact that $\sigma ^{2}_{x} = \sigma ^{2}_{y}=\sigma ^{2}_{z}=1$.
When $k \neq k'$ the operators commute and cancel.
Thus
\[
\sum_{f} \left[B_{fi}(F_{-}) - B_{fi}(F_{+}) \right] = (N_{i}-Z_{i}),       
\]
and
\[
\sum_{f} \left[ B_{fi}(GT_{-}) - B_{fi}(GT_{+}) \right] = 3(N_{i}-Z_{i}).       
\]

The sum-rule for the Fermi matrix elements applies even
when isospin is not conserved.
% --- end paragraph admon ---



% !split
\subsection*{$\beta$-decay}

% --- begin paragraph admon ---
\paragraph{}
For $N > Z$ we usually have
$T_{i}=T_{zi}$ which means that $B(F_{+})=0$.

For $N=Z (T_{zi}=0)$ and $T_{i}=0$ we get 
$B(F_{+})=B(F_{-})=0$, and for $T_{i}=1$ we
have $B(F_{+}) = B(F_{-}) = 2$. Fermi transitions which would be zero
if isospin is conserved are called isospin-forbidden Fermi transitions.

When $N > Z$ there are some situations where one has $B(GT_{+})=0$,
and then we obtain $B(GT_{-}) = 3(N_{i}-Z_{i})$. In particular
for the $\beta_{-}$ decay of the neutron we have $B(F_{-})=1$
and $B(GT_{-})=3$.
% --- end paragraph admon ---










% !split
\subsection*{Core-polarization}

% --- begin paragraph admon ---
\paragraph{}
We need to say something about so-called core-polarization effects.
To do this, we have to introduce elements from many-body perturbation theory.

We assume here that we are only interested in the ground state of the system and 
expand the exact wave function in term of a series of Slater determinants
\[
\vert \Psi_0\rangle = \vert \Phi_0\rangle + \sum_{m=1}^{\infty}C_m\vert \Phi_m\rangle,
\]
where we have assumed that the true ground state is dominated by the 
solution of the unperturbed problem, that is
\[
\hat{H}_0\vert \Phi_0\rangle= W_0\vert \Phi_0\rangle.
\]
The state $\vert \Psi_0\rangle$ is not normalized, rather we have used an intermediate 
normalization $\langle \Phi_0 \vert \Psi_0\rangle=1$ since we have $\langle \Phi_0\vert \Phi_0\rangle=1$.
% --- end paragraph admon ---



% !split
\subsection*{Electromagnetic multipole moments and transitions}

% --- begin paragraph admon ---
\paragraph{}
The Schroedinger equation is
\[
\hat{H}\vert \Psi_0\rangle = E\vert \Psi_0\rangle,
\]
and multiplying the latter from the left with $\langle \Phi_0\vert $ gives
\[
\langle \Phi_0\vert \hat{H}\vert \Psi_0\rangle = E\langle \Phi_0\vert \Psi_0\rangle=E,
\]
and subtracting from this equation
\[
\langle \Psi_0\vert \hat{H}_0\vert \Phi_0\rangle= W_0\langle \Psi_0\vert \Phi_0\rangle=W_0,
\]
and using the fact that the both operators $\hat{H}$ and $\hat{H}_0$ are hermitian 
results in
\[
\Delta E=E-W_0=\langle \Phi_0\vert \hat{H}_I\vert \Psi_0\rangle,
\]
which is an exact result. We call this quantity the correlation energy.
% --- end paragraph admon ---



% !split
\subsection*{Electromagnetic multipole moments and transitions}

% --- begin paragraph admon ---
\paragraph{}
This equation forms the starting point for all perturbative derivations. However,
as it stands it represents nothing but a mere formal rewriting of Schroedinger's equation and is not of much practical use. The exact wave function $\vert \Psi_0\rangle$ is unknown. In order to obtain a perturbative expansion, we need to expand the exact wave function in terms of the interaction $\hat{H}_I$. 

Here we have assumed that our model space defined by the operator $\hat{P}$ is one-dimensional, meaning that
\[
\hat{P}= \vert \Phi_0\rangle \langle \Phi_0\vert ,
\]
and
\[
\hat{Q}=\sum_{m=1}^{\infty}\vert \Phi_m\rangle \langle \Phi_m\vert .
\]
% --- end paragraph admon ---



% !split
\subsection*{Electromagnetic multipole moments and transitions}

% --- begin paragraph admon ---
\paragraph{}
We can thus rewrite the exact wave function as
\[
\vert \Psi_0\rangle= (\hat{P}+\hat{Q})\vert \Psi_0\rangle=\vert \Phi_0\rangle+\hat{Q}\vert \Psi_0\rangle.
\]
Going back to the Schr\"odinger equation, we can rewrite it as, adding and a subtracting a term $\omega \vert \Psi_0\rangle$ as
\[
\left(\omega-\hat{H}_0\right)\vert \Psi_0\rangle=\left(\omega-E+\hat{H}_I\right)\vert \Psi_0\rangle,
\]
where $\omega$ is an energy variable to be specified later.
% --- end paragraph admon ---



% !split
\subsection*{Electromagnetic multipole moments and transitions}

% --- begin paragraph admon ---
\paragraph{}
We assume also that the resolvent of $\left(\omega-\hat{H}_0\right)$ exits, that is
it has an inverse which defined the unperturbed Green's function as
\[
\left(\omega-\hat{H}_0\right)^{-1}=\frac{1}{\left(\omega-\hat{H}_0\right)}.
\]

We can rewrite Schroedinger's equation as
\[
\vert \Psi_0\rangle=\frac{1}{\omega-\hat{H}_0}\left(\omega-E+\hat{H}_I\right)\vert \Psi_0\rangle,
\]
and multiplying from the left with $\hat{Q}$ results in
\[
\hat{Q}\vert \Psi_0\rangle=\frac{\hat{Q}}{\omega-\hat{H}_0}\left(\omega-E+\hat{H}_I\right)\vert \Psi_0\rangle,
\]
which is possible since we have defined the operator $\hat{Q}$ in terms of the eigenfunctions of $\hat{H}$.
% --- end paragraph admon ---



% !split
\subsection*{Electromagnetic multipole moments and transitions}

% --- begin paragraph admon ---
\paragraph{}
These operators commute meaning that
\[
\hat{Q}\frac{1}{\left(\omega-\hat{H}_0\right)}\hat{Q}=\hat{Q}\frac{1}{\left(\omega-\hat{H}_0\right)}=\frac{\hat{Q}}{\left(\omega-\hat{H}_0\right)}.
\]
With these definitions we can in turn define the wave function as 
\[
\vert \Psi_0\rangle=\vert \Phi_0\rangle+\frac{\hat{Q}}{\omega-\hat{H}_0}\left(\omega-E+\hat{H}_I\right)\vert \Psi_0\rangle.
\]
This equation is again nothing but a formal rewrite of Schr\"odinger's equation
and does not represent a practical calculational scheme.  
It is a non-linear equation in two unknown quantities, the energy $E$ and the exact
wave function $\vert \Psi_0\rangle$. We can however start with a guess for $\vert \Psi_0\rangle$ on the right hand side of the last equation.
% --- end paragraph admon ---



% !split
\subsection*{Electromagnetic multipole moments and transitions}

% --- begin paragraph admon ---
\paragraph{}
 The most common choice is to start with the function which is expected to exhibit the largest overlap with the wave function we are searching after, namely $\vert \Phi_0\rangle$. This can again be inserted in the solution for $\vert \Psi_0\rangle$ in an iterative fashion and if we continue along these lines we end up with
\[
\vert \Psi_0\rangle=\sum_{i=0}^{\infty}\left\{\frac{\hat{Q}}{\omega-\hat{H}_0}\left(\omega-E+\hat{H}_I\right)\right\}^i\vert \Phi_0\rangle, 
\]
for the wave function and
\[
\Delta E=\sum_{i=0}^{\infty}\langle \Phi_0\vert \hat{H}_I\left\{\frac{\hat{Q}}{\omega-\hat{H}_0}\left(\omega-E+\hat{H}_I\right)\right\}^i\vert \Phi_0\rangle, 
\]
which is now  a perturbative expansion of the exact energy in terms of the interaction
$\hat{H}_I$ and the unperturbed wave function $\vert \Psi_0\rangle$.
% --- end paragraph admon ---



% !split
\subsection*{Electromagnetic multipole moments and transitions}

% --- begin paragraph admon ---
\paragraph{}
In our equations for $\vert \Psi_0\rangle$ and $\Delta E$ in terms of the unperturbed
solutions $\vert \Phi_i\rangle$  we have still an undetermined parameter $\omega$
and a dependecy on the exact energy $E$. Not much has been gained thus from a practical computational point of view. 

In Brilluoin-Wigner perturbation theory it is customary to set $\omega=E$. This results in the following perturbative expansion for the energy $\Delta E$
\[
\Delta E=\sum_{i=0}^{\infty}\langle \Phi_0\vert \hat{H}_I\left\{\frac{\hat{Q}}{\omega-\hat{H}_0}\left(\omega-E+\hat{H}_I\right)\right\}^i\vert \Phi_0\rangle=
\]
\[
\langle \Phi_0\vert \left(\hat{H}_I+\hat{H}_I\frac{\hat{Q}}{E-\hat{H}_0}\hat{H}_I+
\hat{H}_I\frac{\hat{Q}}{E-\hat{H}_0}\hat{H}_I\frac{\hat{Q}}{E-\hat{H}_0}\hat{H}_I+\dots\right)\vert \Phi_0\rangle. 
\]
% --- end paragraph admon ---



% !split
\subsection*{Electromagnetic multipole moments and transitions}

% --- begin paragraph admon ---
\paragraph{}
\[
\Delta E=\sum_{i=0}^{\infty}\langle \Phi_0\vert \hat{H}_I\left\{\frac{\hat{Q}}{\omega-\hat{H}_0}\left(\omega-E+\hat{H}_I\right)\right\}^i\vert \Phi_0\rangle=\]
\[
\langle \Phi_0\vert \left(\hat{H}_I+\hat{H}_I\frac{\hat{Q}}{E-\hat{H}_0}\hat{H}_I+
\hat{H}_I\frac{\hat{Q}}{E-\hat{H}_0}\hat{H}_I\frac{\hat{Q}}{E-\hat{H}_0}\hat{H}_I+\dots\right)\vert \Phi_0\rangle. 
\]
This expression depends however on the exact energy $E$ and is again not very convenient from a practical point of view. It can obviously be solved iteratively, by starting with a guess for  $E$ and then solve till some kind of self-consistency criterion has been reached. 

Actually, the above expression is nothing but a rewrite again of the full Schr\"odinger equation.
% --- end paragraph admon ---


Defining $e=E-\hat{H}_0$ and recalling that $\hat{H}_0$ commutes with 
$\hat{Q}$ by construction and that $\hat{Q}$ is an idempotent operator
$\hat{Q}^2=\hat{Q}$. 
Using this equation in the above expansion for $\Delta E$ we can write the denominator 
\[
\hat{Q}\frac{1}{\hat{e}-\hat{Q}\hat{H}_I\hat{Q}}=
\]
\[
\hat{Q}\left[\frac{1}{\hat{e}}+\frac{1}{\hat{e}}\hat{Q}\hat{H}_I\hat{Q}
\frac{1}{\hat{e}}+\frac{1}{\hat{e}}\hat{Q}\hat{H}_I\hat{Q}
\frac{1}{\hat{e}}\hat{Q}\hat{H}_I\hat{Q}\frac{1}{\hat{e}}+\dots\right]\hat{Q}.
\]
% !split
\subsection*{Electromagnetic multipole moments and transitions}

% --- begin paragraph admon ---
\paragraph{}
Inserted in the expression for $\Delta E$ leads to 
\[
\Delta E=
\langle \Phi_0\vert \hat{H}_I+\hat{H}_I\hat{Q}\frac{1}{E-\hat{H}_0-\hat{Q}\hat{H}_I\hat{Q}}\hat{Q}\hat{H}_I\vert \Phi_0\rangle. 
\]
In RS perturbation theory we set $\omega = W_0$ and obtain the following expression for the energy difference
\[
\Delta E=\sum_{i=0}^{\infty}\langle \Phi_0\vert \hat{H}_I\left\{\frac{\hat{Q}}{W_0-\hat{H}_0}\left(\hat{H}_I-\Delta E\right)\right\}^i\vert \Phi_0\rangle=
\]
\[
\langle \Phi_0\vert \left(\hat{H}_I+\hat{H}_I\frac{\hat{Q}}{W_0-\hat{H}_0}(\hat{H}_I-\Delta E)+
\hat{H}_I\frac{\hat{Q}}{W_0-\hat{H}_0}(\hat{H}_I-\Delta E)\frac{\hat{Q}}{W_0-\hat{H}_0}(\hat{H}_I-\Delta E)+\dots\right)\vert \Phi_0\rangle.
\]
% --- end paragraph admon ---



% !split
\subsection*{Electromagnetic multipole moments and transitions}

% --- begin paragraph admon ---
\paragraph{}
Recalling that $\hat{Q}$ commutes with $\hat{H_0}$ and since $\Delta E$ is a constant we obtain that
\[
\hat{Q}\Delta E\vert \Phi_0\rangle = \hat{Q}\Delta E\vert \hat{Q}\Phi_0\rangle = 0.
\]
Inserting this results in the expression for the energy results in
\[
\Delta E=\langle \Phi_0\vert \left(\hat{H}_I+\hat{H}_I\frac{\hat{Q}}{W_0-\hat{H}_0}\hat{H}_I+
\hat{H}_I\frac{\hat{Q}}{W_0-\hat{H}_0}(\hat{H}_I-\Delta E)\frac{\hat{Q}}{W_0-\hat{H}_0}\hat{H}_I+\dots\right)\vert \Phi_0\rangle.
\]
% --- end paragraph admon ---



% !split
\subsection*{Electromagnetic multipole moments and transitions}

% --- begin paragraph admon ---
\paragraph{}
We can now this expression in terms of a perturbative expression in terms
of $\hat{H}_I$ where we iterate the last expression in terms of $\Delta E$
\[
\Delta E=\sum_{i=1}^{\infty}\Delta E^{(i)}.
\]
We get the following expression for $\Delta E^{(i)}$
\[
\Delta E^{(1)}=\langle \Phi_0\vert \hat{H}_I\vert \Phi_0\rangle,
\] 
which is just the contribution to first order in perturbation theory,
\[
\Delta E^{(2)}=\langle\Phi_0\vert \hat{H}_I\frac{\hat{Q}}{W_0-\hat{H}_0}\hat{H}_I\vert \Phi_0\rangle, 
\]
which is the contribution to second order.
% --- end paragraph admon ---



% !split
\subsection*{Electromagnetic multipole moments and transitions}

% --- begin paragraph admon ---
\paragraph{}
\[
\Delta E^{(3)}=\langle \Phi_0\vert \hat{H}_I\frac{\hat{Q}}{W_0-\hat{H}_0}\hat{H}_I\frac{\hat{Q}}{W_0-\hat{H}_0}\hat{H}_I\Phi_0\rangle-
\langle\Phi_0\vert \hat{H}_I\frac{\hat{Q}}{W_0-\hat{H}_0}\langle \Phi_0\vert \hat{H}_I\vert \Phi_0\rangle\frac{\hat{Q}}{W_0-\hat{H}_0}\hat{H}_I\vert \Phi_0\rangle,
\]
being the third-order contribution.
% --- end paragraph admon ---





% !split
\subsection*{Learning outcomes}

I hope this is not the case


% inline figure
\centerline{\includegraphics[width=0.6\linewidth]{figures/exam2.jpg}}



% !split
\subsection*{Topics we have covered this year}
\begin{itemize}
\item Single-particle properties and mean-field and relation to data

\item How to set up basis states in second quantization and find expectation values

\item Angular momentum properties and the Wigner-Eckart theorem

\item Short survey of properties of nuclear forces

\item The nuclear shell model

\item And how to relate a shell-model calculation to decays and properties of decays.
\end{itemize}

\noindent
% !split
\subsection*{Final presentation}
\begin{itemize}
\item Introduction with motivation

\item Explain an eventual experimental set up

\item Give a short overview of the theory employed and how it relates to the analysis of eventual data

\item Present and discuss your results

\item Summary, conclusions and perspectives

\item Anything else you think is important. Useful to have backup slides
\end{itemize}

\noindent
In total your talk should have a duration of 20-25 minutes, but longer is also ok. 

% !split
\subsection*{What? Me worry?}


% inline figure
\centerline{\includegraphics[width=0.6\linewidth]{figures/exam1.jpeg}}




% inline figure
\centerline{\includegraphics[width=0.6\linewidth]{figures/whatmeworry.jpeg}}




% !split
\subsection*{Best wishes to you all and thanks so much for your efforts this semester}



% inline figure
\centerline{\includegraphics[width=0.6\linewidth]{figures/Nebbdyr2.png}}


















% ------------------- end of main content ---------------

\end{document}

